\subsection*{S18. Field with differing types across documents}

A path whose observed values have types with no least upper bound, where the
conflicting observations come from different documents.

\begin{lstlisting}[style=json]
{ "id": 1, "value": 42 }
{ "id": 2, "value": "unknown" }
\end{lstlisting}

\options
\begin{enumerate}[label=\textbf{(\alph*)}]
  \item \textbf{Reject}, with an error naming the path and the conflicting
        types. \selected
  \item \textbf{A generated variant} --- \texttt{type value\_ty = Int of int |
        Str of string} --- with the column becoming a tag plus one nullable
        column per alternative.
  \item \textbf{An opaque raw-JSON column.}
  \item \textbf{Widen to string}, serializing non-string values back to JSON
        text.
\end{enumerate}

\decision{Reject, consistently with S16.}

\paragraph{S16 and S18 are one rule.}
Under the path-based framing of this document they are not two scenarios but
one, stated at two kinds of path. All elements of an array share the path
\texttt{.mixed[]}, and all values of a member share the path
\texttt{.value}; in both cases every observation across the whole input feeds
that path, and the requirement is the same:

\begin{quote}
The types observed at a path must have a least upper bound in the type
lattice.
\end{quote}

\noindent
S16 is this rule at an element path, S18 the same rule at a member path.

\paragraph{What does not count as a conflict.}
As in S16, the rejection is narrow. Every mixture below is resolved by a
decision already taken and is accepted:

\begin{itemize}
  \item \texttt{int} and \texttt{float} widen to \texttt{float} (S1a), as does
        an integer literal outside native range (S1b).
  \item Any type together with \texttt{null} or with the key being absent
        becomes an \texttt{option} (S4).
  \item A path never given a type sits below every type and is replaced by any
        observation (S6).
  \item Objects with different members merge into one table whose columns are
        partial (S4).
\end{itemize}

\noindent
What is rejected is a genuine absence of an upper bound: integer against
string, integer against boolean, string against boolean, or any scalar against
an object or an array.

\rationale{The arguments from S16 carry over unchanged. Option (b) would make
the shape of a generated \emph{type} depend on which values happened to
appear, so that adding a document could add a constructor --- the
data-dependence problem of S9 again. Option (d) discards the type information
it claims to preserve, and option (c) leaves the column untyped. A corpus that
does this at a path usually reflects a producer that changed its mind rather
than a deliberate union.}

\limitation{This is the one decision in the document that breaks monotonicity.
Everywhere else, adding documents can only loosen an inferred type: a column
widens from \texttt{int} to \texttt{float}, or from total to partial, and a
schema inferred from a smaller corpus stays valid. Rejection under this rule
can instead invalidate a corpus that worked before, because one new document
can introduce a conflict at a path that was previously uniform.

The scenario is also invisible under single-document input (S0), since it
takes two documents to produce, and the diagnostic must therefore name a pair
of documents and the type each contributed rather than pointing at one place
in one file. Option (b) remains the extension if unions turn out to be worth
supporting, and it would then have to be adopted for S16 at the same time.}
